\section{对边对角模型}

\subsection{无限制的对边对角模型}

\begin{theorem}[无限制的对边对角模型]
    例如：已知 $A = \frac{\pi}{3}, a=3$, 没有其他限制，一般采用余弦定理与基本不等式结合进行求解，对于 $S$ 最大值和 $C$ 的取值范围，使用的基本不等式并不一样。
\end{theorem}

\begin{example}
    已知 $\triangle A B C$ 的内角 $A, ~ B, ~ C$ 的对边分别为 $a, ~ b, ~ c$，且 $A = \frac{\pi}{3}$，$a = 4$，求：
    \begin{itemize}
        \item $b + c$ 的取值范围
        \item $b c$ 的取值范围
        \item 三角形的周长的取值范围
        \item 三角形的面积的取值范围
    \end{itemize}
\end{example}

\v{10em}

\subsection{对边对角模型的动圆性质}

\v{14em}

\begin{example}
    $\triangle A B C$ 的内角 $A, ~ B, ~ C$ 的对边分别为 $a, ~ b, ~ c$ ，已知
    $$
        2 c^2=\left(a^2+c^2-b^2\right)(\tan A+\tan B)
    $$
    \begin{enumerate}
        \item 求 $A$ ；
        \item 若 $\triangle A B C$ 中 $a=2$ ，求 $\triangle A B C$ 面积的取值范围．
    \end{enumerate}
\end{example}

\vspace{12em}

\begin{example}
    已知 $\triangle A B C$ 的内角 $A$、$B$、$C$ 的对边分别为 $a$、$b$、$c$, $a \cos (A+C)=(2 c-b) \cos (B+C)$ ．
    \begin{enumerate}
        \item 求角 $A$ 的大小；
        \item 若 $a=6$ ，求 $\triangle A B C$ 面积的最大值．
    \end{enumerate}
\end{example}

\v{12em}

\subsection{有限制的对边对角模型}

\begin{theorem}[有限制的对边对角模型]
    加上其他限制条件，例如是锐角 $\triangle ABC$，此时只能使用正弦定理求解，最后一般会采用辅助角公式合并成一个 $Asin(\omega x + \varphi)$ 的形式。绘制有效图像分析值域即可。

    可能要求 $\sin A + \sin C$ 或者周长（也就是$a + c$），还可能求 $ac$ 乘积，总之都用正弦定理来做。
\end{theorem}

\begin{example}
    在 $\triangle A B C$ 中，内角 $A, B, C$ 所对的边分别为 $a, ~ b, ~ c, ~ A$ 为锐角，$\triangle A B C$ 的面积为 $S$ ，且
    $$
        \left(b^2+c^2\right) \tan A=4 S+\frac{a^2}{2 \sin 2 A}
    $$
    \begin{enumerate}
        \item 求 $A$ ；
        \item 若 $a=1$ ，求 $S$ 的最大值和周长 $C$ 的取值范围．
    \end{enumerate}
\end{example}

\v{12em}

\begin{example}
    已知 $\overrightarrow{m}=\left(\sin 2 x, \sin ^2 x\right), \overrightarrow{n}=(1,2), f(x)=\overrightarrow{m} \cdot \overrightarrow{n}$ ．
    \begin{enumerate}
        \item 求 $f(x)$ 的最小正周期；
        \item $\triangle A B C$ 中，角 $A$、$B$、$C$ 所对边分别为 $a$、$b$、$c$ ，且 $f(A)=2$ ．
              \begin{enumerate}
                  \item 求角 $A$ 的大小；
                  \item 若 $\triangle A B C$ 为锐角三角形，且 $a=\sqrt{2}$ ，求 $b+\sqrt{2} c$ 的最大值．
              \end{enumerate}
    \end{enumerate}
\end{example}

\v{12em}

\begin{example}
    在 $\triangle A B C$ 中，角 $A, B, C$ 所对的边分别为 $a, b, c$ ，已知 $a \sin B=b \sin 2 A$ ．
    \begin{enumerate}
        \item 求角 $A$ 的大小．
        \item 若 $b=3, ~ \triangle A B C$ 的面积为 $3 \sqrt{3}$ ，求 $\triangle A B C$ 的周长．
        \item 若 $\triangle A B C$ 为锐角三角形，求 $2 \cos B+\cos C$ 的取值范围．
    \end{enumerate}
\end{example}

\v{12em}
